\documentclass[11pt]{article}
\usepackage[letterpaper,landscape,margin=.44in]{geometry}
\usepackage{amsmath,amssymb,xcolor,hyperref,array}
\definecolor{ink}{HTML}{183247}
\definecolor{teal}{HTML}{00685F}
\definecolor{paper}{HTML}{355A82}
\definecolor{wash}{HTML}{EEF4F6}
\definecolor{line}{HTML}{CED9DF}
\hypersetup{colorlinks=true,urlcolor=teal}
\pagestyle{empty}
\setlength{\parindent}{0pt}
\setlength{\parskip}{4pt}
\setlength{\abovedisplayskip}{6pt}
\setlength{\belowdisplayskip}{6pt}
\newcommand{\tr}{\operatorname{tr}}
\newcommand{\scnr}{\operatorname{SCNR}}
\newcommand{\pair}[2]{%
\noindent\begin{minipage}[t]{.475\linewidth}\vspace{0pt}#1\end{minipage}\hfill
\begin{minipage}[t]{.475\linewidth}\vspace{0pt}#2\end{minipage}\par}
\newcommand{\head}[2]{%
{\sffamily\small\bfseries\color{teal} CLUTTER-AWARE ISAC \hfill #1 / 3}\par
{\sffamily\LARGE\bfseries\color{ink} #2}\par\vspace{3pt}
\pair{\colorbox{teal}{\parbox{.95\linewidth}{\color{white}\sffamily\bfseries\strut OUR FIXED-PROPOSAL SOLVER}}}
{\colorbox{paper}{\parbox{.95\linewidth}{\color{white}\sffamily\bfseries\strut LIU ET AL.\hspace{.85em} AO WITH DINKELBACH--SDP}}}\vspace{4pt}}
\newcommand{\sect}[1]{\par{\sffamily\large\bfseries\color{ink} #1}\par\vspace{-3pt}}
\newcommand{\rowline}{\par\vspace{3pt}{\color{line}\hrule}\vspace{5pt}}
\newcommand{\note}[1]{\par\vspace{7pt}\noindent\colorbox{wash}{\parbox{.98\linewidth}{\strut #1\strut}}\par}
\begin{document}
\head{1}{Start with the same sensing objective}
\pair{
\sect{1. Objective and constraints}
Use total transmit covariance $Q_n=W_nW_n^H$ and user covariance $X_{nk}=w_{nk}w_{nk}^H$. For subcarrier $n$, let $a_n,b_n$ be target steering vectors and $u_n$ the receive filter. The objective is
\[
\frac{\sum_n \kappa_n |u_n^Hb_n|^2 a_n^HQ_na_n}
{\sum_n u_n^H[\widehat{\mathcal C}_n(Q_n)+R_{\eta,n}]u_n}.
\]
Here $\kappa_n$ is the target power coefficient. The supplied design map $\widehat{\mathcal C}_n$ maps transmit covariance to cold-clutter covariance. $R_{\eta,n}$ accounts for additive interference and noise.
}{
\sect{1. Objective and constraints}
In the paper, $Q_n\equiv R_{X,n}$ and $X_{nk}\equiv R_{n,k}$. Their sensing SCNR is a ratio of powers summed across subcarriers. They maximize it subject to a user SINR floor and a total power budget.
\[
\mathrm{SINR}_{nk}\ge\gamma_{nk},\hspace{1.6em}
\sum_n\tr(Q_n)\le P.
\]
Here $\gamma_{nk}$ is the required communication SINR and $P$ is the total transmit-power budget. Both methods alternate between updating the receiver and updating the transmitter.

{\small\color{paper}Paper equations (117) and (118).}
}
\rowline
\pair{
\sect{2. Compute the MVDR receiver}
For fixed transmit beams, enforce $u_n^Hb_n=1$. The target term is then fixed. Minimize disturbance power on each tone:
\[
\min_{u_n^Hb_n=1} u_n^H R_{I,n}u_n,\hspace{1.6em}
R_{I,n}=\widehat{\mathcal C}_n(Q_n)+R_{\eta,n}.
\]
Our implementation solves $R_{I,n}v_n=b_n$ and normalizes $v_n$.
}{
\sect{2. Compute the MVDR receiver}
The paper uses the same minimum-variance distortionless response update:
\[
u_n=\frac{R_{I,n}^{-1}b_n}
{b_n^H R_{I,n}^{-1}b_n}.
\]
The normalization preserves the target response. The disturbance covariance includes the clutter generated by the current transmit covariance.

{\small\color{paper}Paper equation (119).}
}
\note{\textbf{Connection.} Changing $Q_n$ changes the cold clutter and therefore the next MVDR filter. Freeze the filters before comparing the two transmit updates.}

\newpage
\head{2}{Convert the transmit problem}
\pair{
\sect{3. Represent covariance using beam powers}
For unit direction $v_{nj}$ with power $p_{nj}\ge0$,
\[
\begin{aligned}
Q_n&=\sum_j p_{nj}v_{nj}v_{nj}^H,\\
X_{nk}&=\sum_{j:o_{nj}=k}p_{nj}v_{nj}v_{nj}^H.
\end{aligned}
\]
The tag $o_{nj}$ assigns a direction to one user or to sensing. With $U$ frozen, the contributions to $F$ and $G$ per watt are $v_{nj}^HA_nv_{nj}$ and $v_{nj}^HC_nv_{nj}$. PSD constraints follow from nonnegative powers. The SINR inequality on the right becomes $Bp\ge d$, with entries
\[
B_{nk,nj}=\big[(1+\gamma_{nk}^{-1})\mathbf1_{o_{nj}=k}-1\big]|h_{nk}^Hv_{nj}|^2.
\]
Entries across tones are zero and $d_{nk}=\sigma_c^2$.
}{
\sect{3. Optimize covariance matrices}
At fixed $U$, write the objective as $F(Q)/G(Q)$, where
\[
\begin{aligned}
F(Q)&=\sum_n\tr(A_nQ_n),\\
G(Q)&=d_0+\sum_n\tr(C_nQ_n).
\end{aligned}
\]
$A_n$ gives target reward, $C_n$ gives cold-clutter cost and $d_0$ is the additive disturbance power after filtering. With $H_{nk}=h_{nk}h_{nk}^H$, rearrange the SINR constraint:
\[
(1+\gamma_{nk}^{-1})\tr(H_{nk}X_{nk})
\ge\tr(H_{nk}Q_n)+\sigma_c^2 .
\]
Optimize $Q_n,X_{nk}\succeq0$ with $Q_n-\sum_kX_{nk}\succeq0$ and the power budget. The SDP drops the user rank-one constraints.

{\small\color{paper}Paper equation (120).}
}
\rowline
\pair{
\sect{4. Charnes--Cooper gives one LP per bank}
The power objective is $a^Tp/(d_0+c^Tp)$, where $a,c$ stack the target and clutter contributions. Set $t=1/(d_0+c^Tp)$ and $y=tp$. Solve
\[
\begin{aligned}
\max_{y,t}\hspace{.85em} & a^Ty\\
\text{s.t.}\hspace{.85em} & By\ge dt,\hspace{.85em} \mathbf1^Ty\le Pt,\\
&d_0t+c^Ty=1,\hspace{.85em} y,t\ge0.
\end{aligned}
\]
Recover powers as $p=y/t$. This derivation uses watts and suppresses the code's numerical scaling and constraint margins.
}{
\sect{4. Dinkelbach gives a sequence of SDPs}
For a trial ratio $\rho$, solve the SDP with objective
\[
F(Q)-\rho G(Q)
=\sum_n\tr[(A_n-\rho C_n)Q_n]-\rho d_0 .
\]
The code omits the additive constant $-\rho d_0$ during the SDP solve and includes $d_0$ when updating the ratio:
\[
\rho \leftarrow \frac{F(Q_{\rm new})}{G(Q_{\rm new})}.
\]
Repeat with the updated $\rho$. At the exact relaxed transmit optimum, the maximum of $F-\rho G$ is zero.

{\small\color{paper}Paper equations (120) and (123).}
}
\note{\textbf{Connection.} Our LP updates the powers for the current candidate directions. We repeat the LP to improve the direction set. The paper repeats the SDP to update the ratio parameter.}

\newpage
\head{3}{Recover beams and continue alternating}
\pair{
\sect{5. Start feasible and expand the bank}
Initialize with zero-forcing user beams and a sensing beam in the user null space, provided the initialization fits within $P$. Add user proposals using the fixed-weight combination of target $A_nW_n$, clutter $-\eta C_nW_n$, and a user-SINR correction. Here $\eta=F/G$ at the current beams.

After solving an LP, its dual prices yield a Hermitian score matrix $M$ for each tone and beam owner. A candidate's best reduced cost is
\[
\max_{\|v\|=1}v^HMv-\lambda=\lambda_{\max}(M)-\lambda,
\]
where $\lambda$ is the power price. A positive reduced cost above the pricing tolerance identifies an eigenvector to add.
}{
\sect{5. Start from covariance and alternate}
The authors' MATLAB baseline initializes from random transmit beams normalized to the total power budget. Each outer round builds the filters from the current covariance and calls the Dinkelbach--SDP transmit solver.

The SDP searches over PSD matrices subject to the covariance constraints. The fixed-weight beam proposal, candidate bank, and LP dual pricing are operations specific to our solver.

{\small\color{paper}Paper equations (119), (120), and (123).}
}
\rowline
\pair{
\sect{6. Recover and check physical beams}
From the LP covariances, use the same user recovery as the paper. Take the dominant residual eigenmodes to form the available sensing streams. Check the recovered beams for SINR, total power, and design SCNR before acceptance. If the check fails, keep the previous beams.

Refresh MVDR after the transmit updates. The six-LP benchmark uses up to three outer rounds with two LPs per round.
}{
\sect{6. Recover physical beams}
Given the optimized covariances, recover one beam per user:
\[
w_{nk}=\frac{X_{nk}h_{nk}}{\sqrt{h_{nk}^HX_{nk}h_{nk}}}.
\]
This preserves that user's desired received power. Recover sensing beams from the dominant eigenmodes of
\[
Q_n-\sum_k w_{nk}w_{nk}^H .
\]
The baseline code recovers the physical beams after its covariance iterations.

{\small\color{paper}Paper equations (121) and (122).}
}
\note{\textbf{What to tell your professor.} The speed comes from solving for scalar powers over selected directions. The LP has a global optimum for its current bank and fixed receiver. Joint transmit--receive global optimality remains unestablished for the finite-budget solver.}
{\small\href{https://arxiv.org/abs/2602.10537}{Liu et al., Section VI-A} \hspace{.85em}
\href{https://github.com/RangLiu0706/Clutter-Aware-ISAC-Tutorial}{Authors' MATLAB benchmark}}
\end{document}
